% !TITLE 鸟鸣涧
% !AUTHOR 王维
% !DYNASTY 唐
% !GENRE 五言绝句
% !ORDER 38

\begin{poem}
人闲桂花落\textsuperscript{1}，夜静春山空。
月出惊山鸟\textsuperscript{2}，时鸣\textsuperscript{3}春涧中。
\end{poem}

\begin{annotations}
\notetext{1}{人闲桂花落：人闲下来才能注意到桂花在飘落。桂花极小，只有在极静的环境中、极闲的心境下，才能感知到花落。}
\notetext{2}{月出惊山鸟：月亮升起的光惊动了山中的睡鸟。“惊”字写月光的明亮和突然。}
\notetext{3}{时鸣：偶尔鸣叫。正因为叫得少，每一声才格外清晰——以动衬静，鸟鸣让山涧显得更寂静。}
\end{annotations}

\begin{appreciation}
这首诗是王维《皇甫岳云溪杂题五首》之一，写于隐居辋川时期。诗题叫鸟鸣涧，但全诗的重点不在鸟鸣，而在寂静。

人闲桂花落，夜静春山空——人闲下来了，桂花轻轻飘落；夜晚寂静，春天的山林显得空旷。桂花落，是一个极轻的动作。只有在极静的环境中，在极闲的心境下，才能听到桂花落地的声音。这不是夸张，是真实的体验：当周围的一切声音都消失时，最细微的动静也会被放大。春山空，空的不只是山，也是人的心。

月出惊山鸟，时鸣春涧中——月亮升起，惊动了山中的鸟儿，它们偶尔在春天的山涧中鸣叫几声。惊字是全诗的关键。月亮怎么会惊动鸟呢？因为月光突然照亮了山林，鸟儿以为是天亮了，或者有什么东西靠近，于是从睡梦中醒来。时鸣，是偶尔鸣叫，不是持续不断地叫。正因为叫得少，每一声才显得格外清晰。

这首诗用了以动衬静的手法。涧中偶尔一声鸟鸣，反而让整个山林显得更加寂静。就像深夜里，一只钟表的滴答声会让房间显得更安静一样。王维是写静的高手，他知道：真正的静，不是万籁俱寂，而是万籁俱寂中的一点微响。

读这首诗，需要放慢呼吸，放慢心跳。想象自己坐在春夜的山涧边，桂花落在肩上，月光照在脸上，远处偶尔传来一声鸟鸣。那一刻，时间仿佛静止了，所有的烦恼都烟消云散。这便是王维想带给我们的：一种超越语言的宁静。
\end{appreciation}
